\documentclass[12pt]{article}

\usepackage[english]{babel}
\usepackage[utf8x]{inputenc}
\usepackage{amsmath}
\usepackage{enumitem}
\usepackage{graphicx}
\usepackage{ulem}
\usepackage{caption}
\usepackage{placeins}
\usepackage[usenames,dvipsnames]{color}
\usepackage[colorinlistoftodos]{todonotes}
\usepackage{listings}
\usepackage{fixltx2e}
\usepackage{scrpage2}
\usepackage{lastpage}
\clearscrheadfoot
\pagestyle{scrheadings}
\usepackage{glossaries}
\usepackage[
top    = 2.75cm,
bottom = 2.00cm,
left   = 2.50cm,
right  = 2.00cm]{geometry}
\setcounter{secnumdepth}{4}


\makeglossaries

\newglossaryentry{glossaryVerweis} {name=abkuerzung, description={Langer Name}}


\begin{document}
\begin{titlepage}
\begin{center}
% Oberer Teil der Titelseite:
\includegraphics[width=0.5\textwidth]{images/logo}\\[1cm]    

\LARGE TGM - HTBLuVA Wien XX \\ IT Department  \\[1.5cm]

% Title
\rule{1.0\textwidth}{1mm}
{ \huge \bfseries \\[0.4cm]  \huge Volkswagen AG \\ \LARGE ERP-Evaluation \\[0.4cm] }

\rule{1.0\textwidth}{1mm}



% Author and supervisor
\noindent 
\vspace{5cm}

\begin{center}
\large
Author: 
Haidn \textsc{Martin} \&
Siegel \textsc{Hannah} \&
Soyka \textsc{Wolfram}
\end{center}

\vfill

% Bottom of the page
{\large \today}

\end{center}
\end{titlepage}

\tableofcontents


%HEADER AND FOOTER
\pagenumbering{arabic}
\ohead{\headmark}
\automark{section}
\ifoot{© Haidn, Siegel, Soyka}
\ofoot{\pagemark ~of \pageref{LastPage}}

\newpage

\section{Auswahl des Unternehmens}

\section{Volkswagen AG}
\subsection{Das Unternehmen}
\subsubsection{Historie}

\subsection{Finanzen}
\subsection{Produktspektrum}
\subsection{Standorte und Unternehmensstruktur}

\subsection{Logistik}
\subsection{Aufbau- und Ablauforganisation}
\subsection{Markt}
%Marktumfeld, -situation, -wert, Trends

\subsection{Ziele und Zukunftsaspekte}

\section{Auswahl eines ERP Systemes}

\section{Projekthandbuch}
\section{Pflichtenheft}
\section{Vorbereitung Kick Off Meeting}




\newpage
\listoftables
\listoffigures
\printglossaries
\subsection{Easy Bibliography}
\begin{thebibliography}{56}

\bibitem{name}
   \textbf{Who}, When\\
  \textit{url}
  \newline last used: dd.mm.yyyy, hh:mm
 
 \bibitem{example} 
  \textbf{Mister Super-genious}, Answer from 20.01.2015\\
  \textit{http://www.stackoverflow.com/question}
  \newline last used: 22.10.2014, 21:00
\end{thebibliography}
\end{document}
